\section{费米子球面张量}\label{sec:fermions}

我们需要一个张量，当若干条腿的位置重合时，可以将这些腿合并。先在射影直线的一个仿射坐标中取标记点 $z_1,\ldots,z_k$。在每个点，输入空间都有一个取半整数值的\emph{能量} $h_i$。各能量决定张量腿的维数，各标记点则决定张量的元素。用于给这些腿定价的度量站点可以独立于 $z_i$ 选取。

证明用到这一构造的两个性质。当一簇标记点发生碰撞时，指定阶数的洛朗系数可经由单独一条腿传递；当其他站点固定时，这条腿的能量由该阶数决定（命题~\ref{prop:fusion}）。反过来，在两点碰撞时，只要让第一个输入在一个有界的能量范围内变化，就可以实现从第二个输入到融合所得张量腿的每一个线性映射。引理~\ref{lem:width} 给出了这个能量界。结合这两个事实，第~\ref{sec:sphere}~节便可把球面张量的价格与移动一条独立连线的价格作比较。

上述半整数分次来自自由费米子~\cite[第~4~节]{kacwang}。Khoruzhii 曾用其模式分解两条连线的交叉 \cite{Khoruzhii}。文献~\cite[定理~\ext{G}{thm:crossing}]{oaG} 中的交叉论证在任意特征下使用了相同的演算。

\subsection{核及其系数}
设 $F$ 为域，取由 $\psi_m$（$m=0,1,\ldots$）生成的外代数。赋予 $\psi_m$ 能量 $m+\tfrac12$，以 $V_h$ 表示能量为 $h\in\frac12\mathbb Z_{\ge0}$ 的子空间。其基由按下标递增次序排列的乘积 $\psi_A=\bigwedge_{m\in A}\psi_m$ 组成，其中 $\sum_{m\in A}(m+\tfrac12)=h$。特别地，$V_0=F$。每个这样的乘积的外次数模二同余于 $2h$。

对 $F$ 中互异的 $z_1,\ldots,z_k$，在 $\prod_i V_{h_i}$ 上定义多线性形式 $C(\mathbf h;\mathbf z)$。对基向量输入，先按站点次序列出所有生成元。其取值是对该列表的所有完美匹配所取的带符号和（威克法则~\cite{wick}）。各因子如下：同一站点的两个生成元的收缩为零；若站点不同，将 $z_i$ 处的 $\psi_m$ 与 $z_j$ 处的 $\psi_n$ 按此次序收缩，所得为
\begin{equation}\label{eq:fermions-kernel}
 [x^m y^n]\,\frac{1}{z_j+y-z_i-x}
 =\frac{(-1)^n\binom{m+n}{m}}{(z_j-z_i)^{m+n+1}} .
\end{equation}
对每个定向匹配，将两两配对的元素相邻排列，取各对收缩之积，再乘以从原列表变到该配对列表的置换的符号。反转一对的方向，收缩与置换符号同时变号。改变各对的先后次序不改变符号。所得的和就是收缩所成的交错矩阵的普法夫式。空列表的值为一，奇数长度列表的值为零。

例如，在四个站点各放一个生成元 $\psi_0$，记 $K_{ij}=(z_j-z_i)^{-1}$。这四个输入的值为
\[
 K_{12}K_{34}-K_{13}K_{24}+K_{14}K_{23}.
\]
当 $z_1$ 与 $z_2$ 碰撞时，第一个匹配含有一个内部配对，从而出现极点，另两个匹配则把该簇与外部相连。这一区别决定了一般情形下的融合展开。

这些公式的系数均为整数，且与基域变换可交换。若一个站点有两个重复的生成元，交换二者的配对伙伴就可将匹配两两抵消。因此即使特征为二，这些多线性形式也诱导出外幂上的形式。交换位于相邻位置的站点 $i,j$ 会产生符号 $(-1)^{(2h_i)(2h_j)}$。把每个系数看成位置的有理函数，则它齐次，次数为 $-\sum_i h_i$，因为一次收缩的次数等于两个能量之和的相反数。仅在站点碰撞处才可能出现极点。作为 $z_i$ 的函数，它在无穷远处有界，且当 $h_i>0$ 时在无穷远处为零。

\subsection{在标记点处融合一个簇}
将站点划分为由连续站点构成的簇 $J_1,\ldots,J_s$，必要时先重排站点。在簇 $J_j$ 中写成 $z_i=p_j+u a_i$，其中中心 $p_j$ 两两不同，每一簇内的偏移量也两两不同。一个偏移量为零的单点簇保持不动。其融合映射在原来的能量上是恒等映射，在其他能量上均为零。

\begin{proposition}[融合]\label{prop:fusion}
对每个簇 $J$，都有一族仅依赖于其偏移量和输入能量的映射
\[
 I_e^J:\bigotimes_{i\in J}V_{h_i}\longrightarrow V_e
 \qquad(e\in\tfrac12\mathbb Z_{\ge0})
\]
置 $H_J=\sum_{i\in J}h_i$。除非 $e-H_J\in\mathbb Z$，否则这些映射为零。在 $u=0$ 处的洛朗展开为
\begin{equation}\label{eq:fermions-fusion}
 C(\mathbf h;\mathbf z)
 =\sum_{\mathbf e}u^{\,\sum_j e_j-\sum_i h_i}\,
 C(\mathbf e;\mathbf p)\circ
       (I^{J_1}_{e_1},\ldots,I^{J_s}_{e_s}).
\end{equation}
对每个系数仅有有限项作出贡献。
\end{proposition}

\begin{proof}
按各簇内部的配对情况将完美匹配分组。固定一个簇 $J$ 及一个这样的部分匹配。在位置 $a_i$ 处，用式~\eqref{eq:fermions-kernel} 收缩其中的各对。把已匹配的对移到最前面，附上置换符号，并保持未匹配生成元的相对次序。在偏移量 $a_i$ 处，对每个未匹配的生成元作替换
\begin{equation}\label{eq:fermions-translation}
 \psi_m\longmapsto
          \sum_{\ell\ge m}\binom{\ell}{m}a_i^{\ell-m}\psi_\ell .
\end{equation}
按原次序对所得向量取外积，仅保留能量 $e$ 的分量。每次平移都可截断到能量不超过 $e$ 的模式，故这是有限维空间上的运算。对一切内部匹配求和，并计入各自带符号的因子，得到 $I_e^J$。若某生成元在同一个输入中重复出现，只要至少有一个副本被收缩，交换两个副本便使符号相反的项成对出现。若二者均未匹配，则其平移后的外积为零。这种消去在任意域上都成立。去掉配对、提高模式下标均不改变能量模 $\mathbb Z$ 的值。

将每个完美匹配分解为内部边和簇间边。在原位置进行一次内部收缩，所得贡献等于上述在偏移量处的收缩乘以 $u$ 的一个幂，其指数是这一对的能量的相反数。在簇间边上，核在 $u=0$ 处正则。关于两个偏移量作泰勒展开时，用到 $[x^m](u a_i+x)^\ell=\binom{\ell}{m}(u a_i)^{\ell-m}$，从而得到式~\eqref{eq:fermions-translation}，且下标每提高一就产生一个 $u$ 因子。因此对于给定的输出能量 $e$，该簇贡献的总指数为 $e-H_J$。将所有簇的内部配对移到前面时，产生的符号都已包含在 $I_e^J$ 中，因为有序对越过任一更早的簇都不改变符号。剩下的符号恰是各输出外积间匹配的符号。这就证明了式~\eqref{eq:fermions-fusion}。在总指数固定时，诸非负数 $e_j$ 的和固定，故每个系数都是有限和。
\end{proof}

例如，把偏移量为 $0,1$ 的两个生成元 $\psi_0,\psi_0$ 融合。它们的总能量为 $1$。内部收缩在能量 $0$ 处产生真空态，系数为 $u^{-1}$。若二者均未匹配，则平移后的外积的第一个非零项为 $\psi_0\wedge\psi_1$，其能量为 $2$，幂次对应于 $u^1$。该项可以与其他中心处的生成元配对。

\subsection{能量与维数}
维数为 $d_h=\dim V_h$ 的一条腿，在 $d_h>0$ 时的宽度是 $\log d_h$。基外积选取若干互异模式，其能量等于这些模式的半整数权重之和。对此类子集计数可知，宽度渐近等于 $c\sqrt h$，其中 $c=\pi/\sqrt3$，对数取自然对数。

\begin{lemma}[能级的维数]\label{lem:dimension}\label{lem:fermions-dimension}
对每个 $h\in\frac12\mathbb Z_{\ge0}$，都有 $d_h\le e^{c\sqrt h}$。当 $h\to\infty$ 遍历所有半整数能级趋于无穷时，有
\[
   d_h\sim
   \frac{\exp(c\sqrt h)}{2^{5/2}3^{1/4}h^{3/4}} .
\]
\end{lemma}

\begin{proof}
对 $t>0$，生成函数为
\[
 Z(t)=\sum_h d_h e^{-th}
     =\prod_{m\ge0}\bigl(1+e^{-t(m+1/2)}\bigr).
\]
函数 $f(x)=\log(1+e^{-x})$ 在 $[0,\infty)$ 上为凸函数。在每个区间 $[mt,(m+1)t]$ 上作中点估计得到 $\log Z(t)\le b/t$，其中 $b=\int_0^\infty f(x)\,dx=\sum_{j\ge1}(-1)^{j+1}/j^2=\pi^2/12$。于是 $d_h\le \exp(th+b/t)$。关于 $t$ 取最小值得到 $h>0$ 时的界，而 $d_0=1$。

将生成元能量加倍，便把能量为 $h$ 的基外积对应到将 $2h$ 分拆成互异奇数部件的分拆。这些分拆与大小为 $2h$ 的自共轭杨图一一对应。每个对角钩的长度为 $2m+1$，其中 $m$ 是它的臂长，严格递减的臂长序列唯一确定杨图。整数 $n$ 的自共轭分拆数的渐近公式的首项为 $\exp(\pi\sqrt{n/6})/(2^{7/4}3^{1/4}n^{3/4})$~\cite[第~5.1~节]{Wildon}。取 $n=2h$ 即得结论。
\end{proof}

\subsection{传输一条腿所需的宽度}
考虑一个两站点簇，偏移量为 $0,1$，输入能量为 $h,g$。把融合到能量 $e$ 的映射记作 $I_e$。固定第一个输入后，此映射成为作用于第二个输入的算子。为在第~\ref{sec:sphere}~节恢复一条连线，我们需要从 $V_g$ 到 $V_e$ 的每一个矩阵单位。下面的引理构造所有这些矩阵单位，且第一个输入在每一能级的宽度至多为 $c(\sqrt e+\sqrt g)+O(1)$。

\begin{lemma}[融合宽度]\label{lem:width}
对 $e,g\in\frac12\mathbb Z_{\ge0}$，置
\[
 H(e,g)=(\sqrt e+\sqrt g)^2+(\sqrt e+\sqrt g)/\sqrt2.
\]
映射
\[
 \bigoplus_{h\le H(e,g)} V_h \longrightarrow
       \operatorname{Hom}(V_g,V_e),\qquad
 (v_h)_h\longmapsto\sum_h I_e(v_h,-)
\]
是满射。取外积基时，每一个矩阵单位都有系数为整数的原像。
\end{lemma}

\begin{proof}
固定模式集 $A,B$，其能量分别为 $e,g$。我们在第一个点构造一个输入，使相应算子把 $\psi_B$ 映到 $\psi_A$，且将 $V_g$ 的其余所有基向量映为零。如果 $e=g=0$，则第一个输入取真空态。其余情形记 $w=\sqrt e$、$v=\sqrt g$。偏移量为零处的输入将取为若干检验向量的外积。用洛朗多项式 $f(X)=\sum_m t_m X^{-m-1}$ 表示检验向量 $\sum_m t_m\psi_m$。在 $I_e$ 的一个矩阵系数中，检验向量有两种可能。它可以与偏移量为一处的输入下标 $b$ 收缩，系数由式~\eqref{eq:fermions-kernel} 给出，为 $[y^b]f(1+y)$。也可以占据输出下标 $\ell$，系数为 $t_\ell$，因为偏移量为零的平移是恒等映射。这一点可由展开输出的外积看出。每个检验向量必须占据这两种位置之一，不同检验向量占据不同的位置。

我们选取检验向量，使每个向量只可使用指定位置或权重严格更大的位置。指定位置恰好用尽可用权重，从而迫使每个检验向量就位。赋予输出下标 $\ell$ 权重 $\tau=(\ell+\tfrac12)/w$，赋予输入下标 $b$ 权重 $\tau=(b+\tfrac12)/v$。某侧能量为零时，该侧没有下标，无需求商。下表给出了每个目标位置对应的检验。
\[
\begin{array}{c|c|c}
\text{目标}& f(X)& r\\ \hline
\ell\in A&X^{-\ell-1}(1-X^{-1})^r&
   \lfloor v\tau-\tfrac12\rfloor+1\\
b\in B&X^{-r-1}(1-X^{-1})^b&
   \lfloor w\tau-\tfrac12\rfloor+1 .
\end{array}
\]
对输出目标，系数 $t_\ell$ 为一，且在 $X=1$ 处的 $r$ 阶零点迫使它能够占据的所有输入位置都具有更大的权重。其他可能的输出下标都超过 $\ell$。对输入目标，有 $f(1+y)=y^b(1+y)^{-r-b-1}$，故其首个输入系数为一。其他可能的输入位置的下标均较大，而 $r$ 的选取保证所有可能的输出位置的权重都更大。这些结论在整数上就成立。

按某个固定次序对所有检验向量取外积。在指定能量的每个矩阵系数中，所有可用位置的权重和为 $w+v$，这也等于所有目标权重之和。按照构造，每个检验向量所占位置的权重至少为其目标权重。由于各位置的权重都为正，每个非零项必须用尽所有位置，并且每个检验向量都必须取到等号。由表可见，即使对侧某个位置具有相同权重，检验向量在此权重下也只能占据其指定位置。特别地，第二个输入的每个生成元均已与一个检验向量收缩。因此，输入集与输出集分别为 $B$ 和 $A$，配对以及外积位置的指派都是唯一的。故所需系数等于 $1$，至多差一个可吸收到所取输入外积中的符号。

由上表中 $r$ 的选取，权重为 $\tau$ 的每个检验向量中，生成元能量至多为 $(w+v)\tau+\tfrac12$。由 $s$ 个生成元构成的外积的能量至少为 $\sum_{m=0}^{s-1}(m+\tfrac12)=s^2/2$，因此检验向量至多有 $\sqrt2(w+v)$ 个。将各自的上界相加，可知输入外积的每个齐次分量的能量都至多为 $H(e,g)$。取这些分量，即得所断言的原像。
\end{proof}

这一截断值中的平方根之和也可通过维数计数得到。如果引理~\ref{lem:width} 的映射在截断值 $H$ 处已是满射，则 $d_e d_g\le \sum_{h\le H}d_h\le (2H+1)e^{c\sqrt H}$。若 $H\le(\sqrt e+\sqrt g)^2$，则取对数并使用引理~\ref{lem:dimension}，在满足 $e+g\to\infty$ 的非零能级上得到 $\sqrt H\ge\sqrt e+\sqrt g-O(\log(e+g+2))$。若 $H$ 更大，同一下界自然成立。引理~\ref{lem:width} 在任意特征下都达到了宽度的这一首项。
